\section{实验结果与分析}

我们把全部定量实验结果和定性示例放在这一章。下面所有结果都是在测试集（总共 1997 条 trial）上完整跑出来的。A2 这边我们取了三个随机种子（42、123、2025）的平均值。

\subsection{实验设置（Experiment Setup）}

我们的软硬件环境列在表~\ref{tab:exp-setup} 里。

\begin{table}[htbp]
  \centering
  \caption{实验软硬件配置}
  \label{tab:exp-setup}
  \begin{tabular}{l l}
    \toprule
    配置项 & 说明 \\
    \midrule
    GPU & NVIDIA RTX 6000 Ada, 48GB 显存 \\
    精度 & bfloat16 混合精度 \\
    CLIP backbone & ViT-B/32（pooled 512 维, visual tokens $50\times512$）\\
    EEG encoder & A2 时间-频谱-空间编码器, 512 维输出 \\
    训练规模 & train 7970, val 1998（全量训练）\\
    测试规模 & test 1997（全量测试）\\
    退化条件 & clean, lowres16, strong\_blur, strong\_noise, occlusion50, mixed \\
    EEG 模式 & vision\_only, real\_eeg, shuffled\_eeg, random\_eeg, eeg\_only \\
    \bottomrule
  \end{tabular}
\end{table}

\subsection{A2 语义识别结果（A2 Semantic Recognition Results）}

A2 在六种退化条件下的语义预测结果我们放在了表~\ref{tab:a2-main}。表里的数值是三个种子平均后的 top-1 accuracy。

\begin{table}[htbp]
  \centering
  \caption{A2 语义融合多种子语义预测结果（三种子均值）}
  \label{tab:a2-main}
  \begin{tabular}{l c c c c c c c}
    \toprule
    退化条件 & Vision-only & Real EEG & Shuffled EEG & Random EEG & Real-Vision & Real-Shuffled & Real-Random \\
    \midrule
    clean & 0.950 & 0.975 & 0.403 & 0.681 & +0.025 & +0.572 & +0.294 \\
    lowres16 & 0.258 & 0.523 & 0.070 & 0.062 & +0.265 & +0.453 & +0.461 \\
    mixed & 0.066 & 0.392 & 0.046 & 0.041 & +0.326 & +0.346 & +0.351 \\
    occlusion50 & 0.811 & 0.880 & 0.237 & 0.378 & +0.069 & +0.643 & +0.502 \\
    strong\_blur & 0.759 & 0.837 & 0.187 & 0.273 & +0.078 & +0.650 & +0.564 \\
    strong\_noise & 0.644 & 0.782 & 0.178 & 0.266 & +0.139 & +0.605 & +0.516 \\
    \bottomrule
  \end{tabular}
\end{table}

从表~\ref{tab:a2-main} 的数据来看，有这么几个现象：(1) real EEG 在所有退化条件下都比 vision-only 高。提升最少的 clean 也有 2.5 个百分点，最多的 mixed 到了 32.6 个百分点——基本上退化越狠，EEG 带来的增益越大；(2) 不管是 shuffled EEG 还是 random EEG，real EEG 在所有条件下都明显比它们高。在 lowres16、mixed、strong\_blur、strong\_noise 这些强退化条件下，两个对照组的准确率远低于 real EEG。clean 条件下对照组靠视觉分支本身就能有不错的分类效果，所以跟 real EEG 的差距没那么大。这说明对照组的性能并不是简单的随机猜——视觉输入还在——只是 EEG 和图像的配对被打乱后，分类能力差不多回到了视觉分支基线。这些对照实验说明，性能提升不是因为多加了一个输入通道或者模型容量变大，而是因为 EEG 和图像确实存在真实的配对关系；(3) 我们做了 3 个种子 $\times$ 6 个条件，总共 18 组评估，real EEG 在所有 18 组里都比任何对照组高。

\placeholderfigure{A2 不同退化条件下 Top-1 accuracy 柱状图}{a2-bar-chart-placeholder}{5.0cm}

\subsection{不同语义融合方法对比（Comparison of Semantic Fusion Methods）}

除了 A2，我们还试了几种不同的融合方法：(1) P2 raw+spectrogram 双路融合——只有时间波形和频谱两路，没有空间分支；(2) P2A2——在 P2 encoder 的基础上加了 A2 风格的三支路再融合；(3) residual adapter——在视觉嵌入上加 EEG 残差适配器，不带门控；(4) prototype bias——在文本原型相似度上加 EEG 预测的偏移项。整体来看，A2 在大部分退化条件下 real-vision gap 最高，而且 real EEG 对 shuffled/random EEG 的优势也最稳定。三支路编码加门控残差融合这套结构设计的好处我们在第 5 章讨论过，实验数据初步说明这个设计方向是对的。

\subsection{Token 级 VTF 结果分析（Token-Level VTF Results）}

VTF 在 token 级别上也做到了 real EEG 高于 vision-only 和 EEG 对照组，这表明 token 级的 EEG 融合是走得通的。不过 VTF 的绝对分类准确率比 A2 的 pooled fusion 低不少（比如 lowres16 条件下 VTF 的 real EEG top-1 大概 0.18，A2 是 0.52）。我们分析可能有两个原因：第一，A2 直接在 CLIP 语义空间里优化类别预测，目标和评价指标一致，token 级融合在小样本情况下训练难度更大，分类稳定性不如 embedding 级的 A2；第二，VTF 只用了一次 cross-attention，A2 用了更深的 gated fusion 网络，信息来回交互更充分。VTF 真正的价值在于它输出的增强视觉 token 可以直接喂给生成式模型，这是给后续生成任务准备的关键中间结构。

\subsection{生成式视觉语言模型结果（Generative VLM Results）}

表~\ref{tab:evlm-results} 是各个生成方案的定量结果。Class-hit 这一列是 strong\_blur 和 strong\_noise 两个强退化条件下的平均 caption 类别命中率。

\begin{table}[htbp]
  \centering
  \caption{生成式视觉语言模型结果对比}
  \label{tab:evlm-results}
  \begin{tabular}{l c c c c c c}
    \toprule
    模型 & 有效率 & 无效率 & Class-hit & Real-Vision & Real-Shuffled & Real-Random \\
    \midrule
    GRU baseline & 68.1\% & 31.9\% & 12.4\% & +2.2\% & +5.2\% & +4.1\% \\
    方案一 Qwen2.5 + LoRA & 32.4\% & 67.6\% & 18.1\% & +2.7\% & +5.3\% & +5.3\% \\
    方案五 full adapter & 91.1\% & 8.9\% & 29.4\% & +5.0\% & +5.5\% & +3.0\% \\
    方案五 LoRA $r=8$ & 91.9\% & 8.1\% & 38.3\% & +9.1\% & +10.3\% & +7.8\% \\
    方案五 LoRA $r=16$ & 83.6\% & 16.4\% & 30.4\% & +5.6\% & +6.0\% & +3.8\% \\
    方案五 多候选重排序 & \textbf{97.8\%} & \textbf{2.2\%} & 37.4\% & +8.4\% & +11.1\% & +10.6\% \\
    \bottomrule
  \end{tabular}
\end{table}

从表~\ref{tab:evlm-results} 能看到几个点：(1) 方案五（Qwen2-VL）比方案一（纯文本 Qwen2.5 + LoRA）好很多。Qwen2-VL 本来就是视觉语言模型，预训练阶段就已经做过视觉-语言对齐；普通的 Qwen2.5 得额外学怎么理解视觉前缀，输出自然不稳定；(2) LoRA $r=8$ 的 caption 类别命中率比 $r=16$ 高。我们猜原因是 $r=8$ 在适配能力和过拟合之间更平衡，$r=16$ 参数太多，数据量小的时候生成稳定性容易掉；(3) 多候选重排序没改模型参数，就把有效率从 91.9\% 提到了 97.8\%，同时类别命中率也没掉太多；(4) 方案五所有变体的 real EEG caption 类别命中率都比 vision-only 和 EEG 对照组高。

这里要说明一下，推理的时候我们没用测试集的真实类别标签。多候选重排序只靠候选 caption 的有效性、重复率、长度，还有跟 A2/VTF 预测的 top-k 语义概念的相似度来排。Caption 类别命中率只在最终评估的时候算，不参与候选选择。

生成式模型的类别命中率还是明显比 A2 的 top-1 准确率低。原因很简单：A2 直接在 CLIP 语义空间里优化分类，目标很明确；生成式模型除了分类还要学语言生成，分类指标低一点很正常。生成式模型的核心价值是能输出自然语言给别人看。

\subsection{定性示例分析（Qualitative Examples）}

表~\ref{tab:qualitative} 里我们放了四个典型的定性例子，直观展示 real EEG 在退化条件下怎么帮上忙的。每个例子都给出了 vision-only、real EEG、shuffled EEG 和 random EEG 四种模式下的 caption。

\begin{table}[htbp]
  \centering
  \caption{定性 caption 示例对比}
  \label{tab:qualitative}
  \begin{tabularx}{\linewidth}{l l X X}
    \toprule
    真实类别 & 退化 & Vision-only caption & Real EEG caption \\
    \midrule
    reflex camera & lowres16 & a hand holding a cell phone & a black and white photo of an old camera \\
    canoe & lowres16 & a plane is flying through the sky & a red and yellow canoe in the water \\
    parachute & mixed & a group of people sitting on steps & a group of parachutes flying in the sky \\
    bolete mushroom & mixed & a man walking down a street at night & a mushroom in the ground among mushrooms \\
    \bottomrule
  \end{tabularx}

  \vspace{4pt}
  \begin{tabularx}{\linewidth}{l l X X}
    \toprule
    真实类别 & 退化 & Shuffled EEG caption & Random EEG caption \\
    \midrule
    reflex camera & lowres16 & a close up of a coffee machine & a close up of a shirt on a hanger \\
    canoe & lowres16 & a boat with boats in the background & a man holding a gun in his hand \\
    parachute & mixed & a group of people sitting around a fire & a woman ironing clothes \\
    bolete mushroom & mixed & a person holding flowers in the air & a girl ironing clothes on a board \\
    \bottomrule
  \end{tabularx}
\end{table}

从表~\ref{tab:qualitative} 能明显看出来，在低分辨率和复合退化条件下，vision-only、shuffled EEG 和 random EEG 生成的 caption 跟真实图像内容差挺多的（比如 canoe 被认成 plane 或 gun，mushroom 被认成 clothes 或 flowers），而 real EEG 的 caption 相对来说更接近正确的物体类别，描述也还算合理。这直观说明了真实配对的 EEG 在退化条件下确实能辅助视觉语义理解。

\subsection{实验结果讨论（Discussion）}

我们把所有实验结果汇总一下，有这么几点：

\textbf{为什么 A2 是定量最强的模型}：A2 直接在 CLIP 语义空间里优化类别预测，目标和评估指标高度一致；生成式模型除了分类还要学语言生成，数据量小的时候很难两个目标都顾好。所以 A2 的分类准确率远超任何生成式模型的 caption 类别命中率。

\textbf{为什么 VTF 定量表现不如 A2}：VTF 设计出来就不是为了纯分类的，而是为了给生成式 VLM 提供合适的输入格式。在小样本情况下，token 级训练比 embedding 级更难收敛，而且 VTF 只有一次 cross-attention，交互深度不如 A2 的门控残差融合。所以 VTF 分类性能不如 A2 是符合预期的。

\textbf{为什么 Qwen2-VL 比 Qwen2.5 好}：Qwen2-VL 预训练阶段就已经做了视觉-语言对齐，视觉 adapter 和语言 decoder 之间通信顺畅。纯文本的 Qwen2.5 没有视觉理解能力，碰到 EEG-enhanced visual prefix 得从头开始学怎么解释视觉信息，生成质量自然不稳。

\textbf{LoRA 秩怎么选}：$r=8$ 在适配能力和过拟合风险之间找到了更好的平衡点；$r=16$ 参数更多，数据量小的时候生成稳定性容易下降，具体表现就是有效输出率变低。

\textbf{训练目标选什么}：用类别名称做目标，caption 类别命中率更高；用类别名称加上过滤后的 BLIP 伪 caption 做目标，文本自然度和有效输出率更好。两个目标各有各的好处，我们分别报告了它们的效果。

\textbf{实验的局限}：所有结果都只在 Thought2Text 的 40 类范围内得到的，不能直接推到其他 EEG 数据集或者更多类别。A2 的多种子评估只用了 3 个种子。生成式模型的分类评估依赖 CLIP 文本原型的零样本分类能力，这个能力受限于 CLIP 在具体类别上的判别水平。

\subsection{实验小结（Summary）}

在我们这个数据集和当前实验设置下：(1) A2 时间-频谱-空间语义融合模型在 18/18 的 seed-condition 组合中 real EEG 都比 vision-only 和 shuffled/random EEG 高，是本次课程设计里最稳的定量模型；(2) EEG 带来的增益和视觉退化程度正相关；(3) 基于 Qwen2-VL 的生成方案在强退化条件下做到了 97.8\% 的有效输出率，real EEG 比 vision-only 的 caption 类别命中率提升了约 8--9 个百分点，但生成式模型在分类指标上还是不如 A2 分类器；(4) 对照实验说明性能提升不是因为多了一个输入通道或者模型变大，而是因为 EEG 和图像之间确实有真实的配对关系。
